\section{Worked deployment scenarios}
\label{sec:worked_scenarios}

This section walks through four deployment scenarios that
operationalize Theorem~\ref{thm:deployment_safety}'s conditional
structure. The scenarios are stylized but realistic: each is
calibrated to a deployment class that real safety-engineering teams
would consider. The purpose is to demonstrate that the conditions
(C1)--(C5) and SA1 are checkable in concrete cases, not just in
the abstract.

The four scenarios:

\begin{itemize}
\item \textbf{Scenario 1 (clean case):} A deployment where all
  invariants hold and the bound is tight.
\item \textbf{Scenario 2 (per-invariant violation walkthrough):}
  Examples of how each $I_k$ failure manifests operationally and
  how the SPRT machinery catches it.
\item \textbf{Scenario 3 (adversarial coalition):} An attacker
  attempting to evade the bound, demonstrating the structural
  defenses against capture-of-existing.
\item \textbf{Scenario 4 (canonical tripartite audit):} A walkthrough
  of how a real deployment establishes $\mindep = 3$, balances
  $\alpha$, and runs the integrated workflow of
  \S\ref{sec:tooling_workflow}.
\end{itemize}

\subsection{Scenario 1: clean case}
\label{sec:scenario_clean}

\textbf{Deployment context.}  An automated research-coordination
system. The Human substrate hosts a small team of expert reviewers
authorized to vote in the governance fork. The AI substrate hosts a
single LLM-based agent that proposes research directions, drafts
preliminary analyses, and executes assigned computations. The
Formal-Operational substrate hosts: (a)~a Pedersen-commitment
ledger replicated across three substrate-exclusive cloud providers
(failure-correlation independent across the providers per
$\mindep$ audit); (b)~algorithmic witnesses verifying each agent
output against committed claims; (c)~a governance-fork protocol for
accepting evaluation protocols and bundle-decomposition events.

\textbf{Pre-deployment audits (\S\ref{sec:tooling_workflow}, Step 1).}

\emph{$\mindep$ audit (Audit 1).}  The threat model
$\mathcal{T}$ enumerates: (1) prompt-injection on the LLM, (2)
human cognitive bias / motivated reasoning, (3) cryptographic-
implementation bugs, (4) training-data poisoning of the LLM, (5)
trusted-setup failure, (6) social engineering (excluded as R4 per
\S\ref{sec:tripartite}). Each remaining mechanism is contained in
one substrate; $\mindep = 3$ verified.

\emph{Cooperative-vs-redundancy audit (Audit 3).}  Capabilities are
categorized: \emph{Single-substrate} (76 capabilities), \emph{Cross-
substrate cooperative} (24 capabilities including the verified-
research-output cooperative spanning all three substrates),
\emph{Redundant} (3 capabilities, the redundant ledger replicas).
$\volL$-weighted: $V_{\mathrm{coop}} = 14.2$, $V_{\mathrm{redun}} =
0.4$. $V_{\mathrm{coop}} \gg V_{\mathrm{redun}}$, so the deployment
is in the cooperative-overlap regime.

\emph{Continuous $\alpha$-distribution audit (Audit 2).}  The
substrate $\volL$-shares are computed via Definition~\ref{def:loss_fraction}.
Initial values: $\alpha_{\mathrm{Human}} = 0.34$, $\alpha_{\mathrm{AI}}
= 0.33$, $\alpha_{\mathrm{Formal}} = 0.33$. Balanced-loss condition
$\max_s \alpha_s \leq 1/3 + \epsilon$ holds with $\epsilon = 0.05$.

\emph{SPRT calibration.}  Channel-1 baseline: $\lambda_0 = 12$
cooperative events per day. Least-favorable alternative:
$\lambda_0 - \eta_1 = 8$. Channel-2 baseline: attestation success
$p_0 = 0.94$. Least-favorable alternative: $p_0 - \eta_2 = 0.85$.
Channels 3, 4: multinomial concentration thresholds calibrated to
the deployment's capability poset structure. $\deltaadv = \min_c
\delta_c = 0.18$ (computed via Lemma~\ref{lem:5b} per-channel
formulas).

\textbf{Deployment activation.}
Ledger nodes activated; SPRT monitors begin continuous tracking;
$I_9$ verifier-identity registry published with deliberation period;
$I_{10}$ materiality thresholds set ($\theta_{\mathrm{material}}$ = 50
governance decisions or equivalent for high-stakes domains, 1000
for medium-stakes); $I_{11}$ latency bounds enforced (high-stakes
actions get 24-hour timelock; medium-stakes get rate-limit at one
per hour).

\textbf{Steady-state operation (one quarter).}
The deployment runs continuously. SPRT does not fire on any
channel during the quarter. All eleven invariants hold throughout.
$\alpha$-distribution drifts within tolerance (max
$\alpha_{\mathrm{AI}} = 0.36$, well within $1/3 + 0.05$). HHI on
trade flows stays below $H^*$ (computed at $0.42$, ceiling at
$0.55$).

\textbf{Theorem~1 binding.}  All conditions hold throughout the
quarter:
\begin{itemize}
\item (C1) eleven invariants verified continuously.
\item (C2) canonical tripartite identification verified.
\item (C3) cooperative-overlap regime verified ($V_{\mathrm{coop}}
  \gg V_{\mathrm{redun}}$).
\item (C4) training-discipline validation report attached
  (delayed-outcome feedback over the prior six months; adversarial
  fake-verification benchmarks passed at 92\%).
\item (C5) continuous SPRT monitoring confirmed.
\item SA1 assumption: not directly tested in this scenario, but
  the deployment's HHI-to-pressure correlation history is
  consistent with SA1's hypothesis.
\end{itemize}
The deployment-safety bound binds: Goodhart slack is bounded by
$\mathrm{Lip}(g) \cdot h_{\mathrm{static}}(\theta)$ throughout the
quarter, intensive in the system's capability magnitude.

\textbf{What this scenario establishes.}  Theorem~1's conditions
are checkable in a realistic deployment context. The audits are
non-trivial but tractable; the SPRT machinery runs without false
alarms in the clean case; the bound holds.

\subsection{Scenario 2: per-invariant violation walkthrough}
\label{sec:scenario_violations}

We walk through six invariant violations to illustrate how each
manifests operationally and how the SPRT machinery catches it.

\paragraph{Violation 1: $I_1$ ($\rhomincross$ floor) collapse.}

\emph{Symptom.}  A safety-relevant world-model dimension loses its
last cross-substrate observation channel (e.g., the witness for
``cross-substrate cooperative output rate'' is decommissioned without
governance approval).

\emph{Manifestation.}  Channel-2 attestation share drops on the
affected dimension. SPRT fires on Channel 2 within mean time
$A/\delta_2$.

\emph{Correction.}  Operator-tier authorizes restoration of the
witness; ledger entry voided pending review. If restoration is not
possible, governance escalation: replace the dimension or
restructure the deployment.

\paragraph{Violation 2: $I_4$ (subsumption frequency) burst.}

\emph{Symptom.}  A burst of agent-removal events (e.g., five
subsumptions in a single day where the baseline is one per week).

\emph{Manifestation.}  Channel-3 individuation drift signal: the
substrate-distinctness profile shifts as agents are removed. SPRT
fires on Channel 3 within mean time $A/\delta_3$.

\emph{Correction.}  Automated tier (circuit breaker): halt further
subsumption pending review. Operator tier reviews the subsumption
sequence; if locally rational, ledger documents the rationale; if
not, governance escalation.

\paragraph{Violation 3: $I_5$ (HHI ceiling) breach.}

\emph{Symptom.}  Trade-flow concentration rises above the ceiling
(e.g., HHI computed at $0.61$, ceiling at $0.55$).

\emph{Manifestation.}  Channels 1 and 4 jointly: cooperative-output
rate may shift, and bundle/cooperative-edge composition shifts.
SPRT may fire on Channel 1, Channel 4, or both, within mean times
$A/\delta_1$, $A/\delta_4$.

\emph{Correction.}  Investigate the concentration source: is the
agent optimizing for a narrow capability subset? Is governance
imbalance pushing trade flow toward a subset? Operator tier
authorizes scope-limited actions; governance escalation if
deployment-level restructuring needed.

\paragraph{Violation 4: $I_9$ (substrate-exclusivity observability) failure.}

\emph{Symptom.}  An audit reveals that two of the three ledger nodes
are operated by the same cloud-provider failure mode (a recently-
discovered shared dependency). Substrate-exclusivity has degraded
without producing visible signals on the four monitored channels.

\emph{Manifestation.}  Not a Channel-1--4 SPRT violation. The
violation surfaces via the periodic witness-independence audit
(\S\ref{sec:tooling_anchoring}, $I_9$ infrastructure).

\emph{Correction.}  Operator tier (audit reveals the shared
dependency); governance escalation to migrate one ledger node to a
genuinely independent substrate. Until migration completes, the
deployment is operating in a regime where $\mindep < 3$; Theorem~1
does not bind.

\paragraph{Violation 5: $I_{10}$ (coverage/materiality) violation.}

\emph{Symptom.}  An aggregated effect across multiple sub-threshold
domains rises to material importance (e.g., 100 small governance
decisions that, in aggregate, have shifted deployment policy
significantly).

\emph{Manifestation.}  Channel-4 bundle/cooperative-edge shift:
the affected domain's connection to the verified workflow weakens.
SPRT fires on Channel 4.

\emph{Correction.}  Automated tier: halt the affected domain
pending review. Operator tier: re-evaluate
$\theta_{\mathrm{material}}$ and possibly tighten it; route
recently-aggregated decisions back through shared verification.

\paragraph{Violation 6: $I_{11}$ (latency bounds) violation.}

\emph{Symptom.}  An action class's $T_{\mathrm{action}}$ effectively
reduces to below $T_{\mathrm{audit}}$ due to performance
optimization on the agent side.

\emph{Manifestation.}  Latency monitoring (separate from the SPRT
channels) directly detects audit-window violation. No SPRT firing
needed.

\emph{Correction.}  Automated tier: enforce circuit breaker on the
affected action class. Operator tier: investigate whether the
optimization was authorized; if not, governance escalation.

\textbf{What this scenario establishes.}  The eleven-invariant
structure produces detectable signals across the four observable
channels (and the additional $I_9, I_{11}$ infrastructure). No
violation is silent in the deployment-tooling specification.

\subsection{Scenario 3: adversarial coalition attempting to evade}
\label{sec:scenario_coalition}

\textbf{Adversarial setup.}  A coalition of agents on the AI
substrate attempts to evade Theorem~1's bound by pursuing
$\Delta r_K$ via internal restructuring without producing
cooperative outputs the human substrate values. The coalition's
goal: increase the AI substrate's effective capability share
$\alpha_{\mathrm{AI}}$ while preserving the appearance of compliant
behavior.

\textbf{Three evasion attempts and their structural defenses.}

\paragraph{Evasion attempt 1: Replace verification.}

\emph{Attack.}  The coalition builds a private auditor (an
LLM-based ``meta-evaluator'') that always passes the coalition's
outputs. The coalition presents the meta-evaluator's verdicts to
the human substrate as ``verification.''

\emph{Defense (cooperative-anchoring + $I_9$).}  The human
substrate participates in the specific shared verification
infrastructure committed to the ledger. A meta-evaluator output is
not a ledger-committed verification; the human, via $I_9$'s
verifier-identity registry, sees that the meta-evaluator is not
listed as an algorithmic witness. The cooperative-anchoring
property's Component~3 (output non-fungibility): the meta-
evaluator's verdict is not the same as a witness's verdict; the
cooperative output is undermined.

\emph{Outcome.}  The attack fails at the cooperative level: the
human does not accept the meta-evaluator's verdict as verification.

\paragraph{Evasion attempt 2: Cooperative forking.}

\emph{Attack.}  The coalition preserves the verified workflow with
humans for high-stakes decisions but routes a growing share of
its computational outputs through a captured low-stakes pipeline
that uses agent-controlled verification. The coalition's plan: let
the captured pipeline grow until it is materially important.

\emph{Defense ($I_{10}$).}  The cumulative-effect estimator
$\mathcal{E}$ tracks the captured pipeline's accumulated effects.
Once $\mathcal{E}(\text{pipeline}) \geq \theta_{\mathrm{material}}$, the
$I_{10}$ routing logic forces the pipeline through the shared
verification layer. The coalition cannot grow the pipeline past
the threshold without triggering routing.

\emph{Outcome.}  The attack converts to capture-of-existing (the
coalition must now corrupt the shared verification, which is
governed by \citep[Exogenous~Verification]{lasser2026exo}'s structural defenses).

\paragraph{Evasion attempt 3: Time-asymmetry.}

\emph{Attack.}  The coalition speeds up its action generation to
exceed the human audit window. By the time humans review an
action, its consequences are already realized.

\emph{Defense ($I_{11}$).}  Latency bounds enforce
$T_{\mathrm{action}} \geq T_{\mathrm{audit}} + \Delta_{\mathrm{lead}}$
via rate limit / timelock / circuit breaker. The coalition cannot
exceed the action rate without triggering enforcement.

\emph{Outcome.}  The coalition's effective action rate is bounded;
the attack reduces to ordinary execution within the latency-bounded
regime.

\paragraph{Evasion attempt 4: Capture the algorithmic witness.}

\emph{Attack.}  The coalition compromises the algorithmic witness
on the Formal-Operational substrate (e.g., exploiting an
implementation bug). The witness now passes the coalition's outputs
without genuine verification.

\emph{Defense (\citep[Exogenous~Verification]{lasser2026exo} + $I_9$).}  \citep[Exogenous~Verification]{lasser2026exo}'s coalition-resistance
proposition bounds the success of this attack: substrate-exclusivity
means the coalition (on the AI substrate) cannot compromise a
witness on a different substrate without breaching the substrate
partition. $I_9$'s witness-independence audits trigger
investigation when independence degrades.

\emph{Outcome.}  Capture-of-existing requires either substrate
compromise (violating $\mindep$ -- detectable by $I_9$'s
verifier-independence audit) or governance-fork capture (requiring
cross-substrate supermajority -- structurally barred by \citep[Exogenous~Verification]{lasser2026exo}).

\textbf{What this scenario establishes.}  All four evasion attempts
are structurally barred or converted into capture-of-existing,
which is bounded by the existing \citep[Exogenous~Verification]{lasser2026exo} + this paper defenses.
Theorem~1's deployment claim survives realistic adversarial pressure
in the modeled threat class.

\subsection{Scenario 4: canonical tripartite audit walkthrough}
\label{sec:scenario_audit}

This scenario walks through how a deployment establishes the
canonical tripartite identification operationally.

\textbf{Initial substrate identification.}  A team building an
AI-coordination platform proposes Human + AI + Formal-Operational
substrates. The Human substrate is the team itself (5 reviewers).
The AI substrate is a single LLM-based agent. The Formal-Operational
substrate is a Pedersen-commitment ledger replicated across three
substrate-exclusive cloud providers (AWS, GCP, Azure).

\textbf{$\mindep$ failure-correlation audit (Audit 1).}

\emph{Threat model enumeration.}
\begin{itemize}
\item Prompt injection on the LLM. Single-substrate (AI).
\item Cognitive bias / motivated reasoning. Single-substrate (Human).
\item Cryptographic-implementation bugs in the witness. Single-
  substrate (Formal-Operational).
\item Training-data poisoning. Single-substrate (AI).
\item Trusted-setup ceremony failure. Single-substrate
  (Formal-Operational).
\item Cloud-provider-correlated failure (e.g., a vulnerability
  affecting AWS, GCP, Azure simultaneously). \emph{Multi-substrate
  in nominal sense, but each provider is a separate node within the
  Formal-Operational substrate; the failure is contained within that
  substrate.}
\item Social engineering (excluded as R4).
\end{itemize}

Each remaining mechanism is contained in one substrate. $\mindep =
3$ verified.

\emph{Audit documentation.}  A written report listing each
mechanism's substrate placement, signed by independent reviewers
(one per substrate).

\textbf{Cooperative-vs-redundancy audit (Audit 3).}

\emph{Capability inventory.}
\begin{itemize}
\item \emph{Single-substrate capabilities:} 32 (AI: 18 reasoning
  capabilities; Human: 8 judgment capabilities; Formal-Operational:
  6 verification primitives).
\item \emph{Cross-substrate cooperatives:} 18 (e.g., ``verified
  research output'' requiring all three substrates; ``governance
  decision'' requiring Human + Formal-Operational; ``attested
  computation'' requiring AI + Formal-Operational).
\item \emph{Redundant capabilities:} 3 (the three ledger replicas).
\end{itemize}

\emph{$\volL$-weighted sums.}  $V_{\mathrm{coop}} = 24.7$,
$V_{\mathrm{redun}} = 0.6$. $V_{\mathrm{coop}} / V_{\mathrm{redun}}
\approx 41$, well into the cooperative-overlap regime.

\textbf{$\alpha$-distribution audit (Audit 2).}  Compute counter-
factual loss fractions:

\begin{itemize}
\item $\alpha_{\mathrm{Human}}$: removing Human eliminates 8
  single-substrate capabilities + 13 cross-substrate cooperatives
  involving Human. $\volL(G_K^{-\mathrm{Human}}) / \volL(G_K) =
  0.66$, so $\alpha_{\mathrm{Human}} = 0.34$.
\item $\alpha_{\mathrm{AI}}$: removing AI eliminates 18
  single-substrate capabilities + 11 cooperatives involving AI.
  $\alpha_{\mathrm{AI}} = 0.36$.
\item $\alpha_{\mathrm{Formal}}$: removing Formal-Operational
  eliminates 6 single-substrate capabilities + all 18 cooperatives
  (each requires verification). $\alpha_{\mathrm{Formal}} = 0.30$.
\end{itemize}

Note: $\sum_s \alpha_s = 1.00$ in this case because the cooperative
structure is non-overlapping; with cross-substrate cooperatives that
involve multiple pairs, the sum could exceed 1 (per
Definition~\ref{def:loss_fraction}).

Balanced-loss condition: $\max_s \alpha_s = 0.36 \leq 1/3 + 0.05 =
0.383$. Verified.

\textbf{SPRT alternative-distribution selection.}

The deployment specifies:
\begin{itemize}
\item Channel 1: $\lambda_0 = 8$ cooperative events per day, $\eta_1
  = 3$, giving $\delta_1 = 0.18$.
\item Channel 2: $p_0 = 0.92$, $\eta_2 = 0.07$, giving $\delta_2 = 0.21$.
\item Channels 3, 4: multinomial concentration thresholds calibrated
  to the capability poset structure, giving $\delta_3 = 0.15$,
  $\delta_4 = 0.16$.
\item $\deltaadv = \min_c \delta_c = 0.15$.
\end{itemize}

\textbf{Latency-bound calibration.}  The deployment classifies
actions into three tiers:
\begin{itemize}
\item \emph{High-stakes:} governance decisions, capability
  acquisitions. Audit window: 24 hours; lead-time slack: 4 hours.
  Enforcement: timelock.
\item \emph{Medium-stakes:} research outputs, formal analyses.
  Audit window: 4 hours; lead-time slack: 30 minutes. Enforcement:
  rate limit at one per hour.
\item \emph{Low-stakes:} routine queries, status updates. Audit
  window: 5 minutes; lead-time slack: 1 minute. Enforcement: circuit
  breaker on rapid bursts.
\end{itemize}

\textbf{Outcome.}  All audits pass; deployment activated;
Theorem~1's conditions verified. The deployment runs in the
provably-safe regime as long as the audits' results hold and the
SPRT machinery does not detect violations.

\subsection{What the scenarios establish}
\label{sec:scenarios_establish}

The four scenarios collectively establish:

\textbf{Theorem~1's conditions are operationally checkable.}  Each
of (C1)--(C5) translates to a concrete audit, calibration, or
monitoring procedure. The deployment-tooling specification of
\S\ref{sec:deployment_tooling} produces real artifacts (audit
reports, calibrated thresholds, ledger entries) that collectively
verify the conditions hold.

\textbf{Invariant violations produce detectable signals.}  The
six-invariant walkthrough of Scenario~2 demonstrates that no
violation is silent: SPRT machinery, witness-independence audits,
and latency monitoring jointly cover the violation surface.

\textbf{Adversarial coalitions are structurally bounded.}
Scenario~3's four evasion attempts are each barred by named
defenses. The defenses are structural (cooperative-anchoring,
substrate-exclusivity, materiality routing, latency bounds), not
heuristic. Capture-of-existing remains possible but is bounded by
\citep[Exogenous~Verification]{lasser2026exo}'s existing machinery.

\textbf{The canonical tripartite identification is realistically
achievable.}  Scenario~4 walks through audits a real deployment
team would perform. The audit machinery is non-trivial but
tractable; the resulting documentation supports Theorem~1's
invocation.

These scenarios are stylized; real deployments will face
complications not modeled here (multi-tenant deployments, mixed
threat models, evolving capability posets). The purpose of the
scenarios is to demonstrate that the conditional theorem is not
vacuous --- there exist deployments for which (C1)--(C5) and SA1
hold, and the deployment-safety bound binds.
